\documentclass[10pt,letterpaper]{article}

% ─── Packages (XeTeX compatible) ──────────────────────────────────────────────
\usepackage{fontspec}
\usepackage{xunicode}
\usepackage{mathtools}
\usepackage{amssymb}
\usepackage[table,xcdraw]{xcolor}
\usepackage{geometry}
\usepackage{graphicx}
\usepackage{hyperref}
\usepackage{listings}
\usepackage{booktabs}
\usepackage{array}
\usepackage{tabularx}
\usepackage{enumitem}
\usepackage{fancyhdr}
\usepackage{titlesec}
\usepackage{tcolorbox}
\usepackage{tikz}
\usepackage{mdframed}
\usepackage{float}
\usepackage{setspace}
\usepackage{microtype}
\usepackage{parskip}
\usepackage{caption}
\usetikzlibrary{shapes.geometric, arrows.meta, positioning, fit, backgrounds, calc}

% ─── Colour Palette — pure B&W for text, colours only in TikZ diagram ─────────
\definecolor{Black}{HTML}{000000}
\definecolor{DarkGray}{HTML}{222222}
\definecolor{MidGray}{HTML}{555555}
\definecolor{LightGray}{HTML}{999999}
\definecolor{RuleGray}{HTML}{CCCCCC}

% Diagram-only colours
\definecolor{DiagFE}{HTML}{2563EB}       % frontend blue
\definecolor{DiagFELight}{HTML}{DBEAFE}
\definecolor{DiagIPC}{HTML}{0891B2}      % IPC teal
\definecolor{DiagIPCLight}{HTML}{CFFAFE}
\definecolor{DiagBE}{HTML}{EA580C}       % backend rust
\definecolor{DiagBELight}{HTML}{FFEDD5}
\definecolor{DiagDB}{HTML}{16A34A}       % db green
\definecolor{DiagDBLight}{HTML}{DCFCE7}
\definecolor{DiagProc}{HTML}{7C3AED}     % subprocess purple
\definecolor{DiagProcLight}{HTML}{EDE9FE}
\definecolor{DiagAxum}{HTML}{D97706}     % axum amber
\definecolor{DiagAxumLight}{HTML}{FEF3C7}

% ─── Hyperref ───────────────────────────────────────────────────────────────────
\hypersetup{
  colorlinks=false,
  pdfborder={0 0 0},
  pdftitle={Synclime: Technical Architecture Report},
  pdfauthor={Ahmed Trooper},
  pdfsubject={Software Architecture and Systems Engineering},
}

% ─── Code Style — pure B&W ──────────────────────────────────────────────────────
\lstdefinestyle{code}{
  backgroundcolor=\color{white},
  basicstyle=\ttfamily\footnotesize\color{Black},
  keywordstyle=\bfseries\color{Black},
  commentstyle=\itshape\color{MidGray},
  stringstyle=\color{DarkGray},
  numberstyle=\tiny\color{LightGray},
  numbers=left,
  numbersep=8pt,
  frame=single,
  framerule=0.4pt,
  rulecolor=\color{RuleGray},
  breaklines=true,
  breakatwhitespace=true,
  tabsize=4,
  showstringspaces=false,
  captionpos=b,
}
\lstset{style=code}

% ─── Geometry ───────────────────────────────────────────────────────────────────
\geometry{margin=1.2in, top=0.8in, bottom=0.8in}

% ─── Header / Footer ────────────────────────────────────────────────────────────
\pagestyle{fancy}
\fancyhf{}
\fancyhead[L]{\small\color{MidGray}\sffamily\textbf{Synclime} \textemdash{} Technical Architecture}
\fancyhead[R]{\small\color{MidGray}\sffamily v1.0.0}
\fancyfoot[C]{\small\color{MidGray}\thepage}
\renewcommand{\headrulewidth}{0.5pt}
\renewcommand{\headrule}{\hbox to\headwidth{\color{RuleGray}\leaders\hrule height \headrulewidth\hfill}}

% ─── Section Styling ────────────────────────────────────────────────────────────
\titleformat{\section}
  {\LARGE\bfseries\sffamily\color{Black}}
  {\thesection.}{0.5em}{}[\vspace{2pt}\color{RuleGray}\titlerule]
\titleformat{\subsection}
  {\large\bfseries\sffamily\color{DarkGray}}
  {\thesubsection.}{0.5em}{}
\titleformat{\subsubsection}
  {\normalsize\bfseries\sffamily\color{MidGray}}
  {\thesubsubsection.}{0.4em}{}

% ─── Callout box — white bg, black border, B&W ──────────────────────────────────
\tcbuselibrary{breakable, skins}

\newtcolorbox{insightbox}[1]{
  enhanced, breakable,
  colback=white,
  colframe=Black,
  fonttitle=\bfseries\sffamily\color{Black},
  title={#1},
  coltitle=Black,
  boxrule=0.8pt,
  arc=0pt,
  outer arc=0pt,
  left=10pt, right=10pt, top=6pt, bottom=6pt,
  attach boxed title to top left={yshift=-2mm, xshift=6mm},
  boxed title style={colback=white, colframe=Black, arc=0pt, outer arc=0pt, boxrule=0.8pt}
}
\newtcolorbox{techbox}[1]{
  enhanced, breakable,
  colback=white,
  colframe=MidGray,
  fonttitle=\bfseries\sffamily\color{Black},
  coltitle=Black,
  title={#1},
  boxrule=0.4pt,
  arc=0pt,
  left=10pt, right=10pt, top=6pt, bottom=6pt,
  attach boxed title to top left={yshift=-2mm, xshift=6mm},
  boxed title style={colback=white, colframe=MidGray, arc=0pt, outer arc=0pt, boxrule=0.4pt}
}

% ═══════════════════════════════════════════════════════════════════════════════
\begin{document}
\setstretch{1.2}

% ─── TITLE PAGE — pure white / black ────────────────────────────────────────────
\begin{titlepage}
  \color{Black}
  \vspace*{3cm}
  \begin{center}
    {\fontsize{11}{14}\selectfont\sffamily\color{MidGray}\MakeUppercase{Technical Architecture Report}}
    \vspace{0.6cm}

    \noindent{\color{Black}\rule{0.5\linewidth}{1pt}}
    \vspace{0.7cm}

    {\fontsize{44}{52}\selectfont\bfseries\sffamily SYNCLIME}

    \vspace{0.5cm}
    {\fontsize{13}{17}\selectfont\color{DarkGray}
      A Cross-Platform Native Desktop Media Extraction Engine\\[3pt]
      Built on Tauri v2, SolidJS, Rust 2021 Edition, and SQLite
    }

    \vspace{0.7cm}
    \noindent{\color{Black}\rule{0.5\linewidth}{1pt}}
    \vspace{1.6cm}

    \begin{tabular}{rl}
      \color{MidGray}\textbf{Author}    & Ahmed Trooper \\[5pt]
      \color{MidGray}\textbf{Version}   & 1.0.0 \\[5pt]
      \color{MidGray}\textbf{Date}      & July 2026 \\[5pt]
      \color{MidGray}\textbf{Repository}& github.com/AhmedTrooper/Synclime \\[5pt]
      \color{MidGray}\textbf{License}   & MIT / Apache 2.0 (Dual) \\[5pt]
      \color{MidGray}\textbf{Platforms} & Windows x64/ARM64 \textbullet{} macOS ARM64 \textbullet{} Linux x64/ARM64 \\
    \end{tabular}

    \vspace{1.8cm}
    \begin{mdframed}[linecolor=Black, linewidth=0.8pt, innerleftmargin=14pt,
                    innerrightmargin=14pt, innertopmargin=10pt, innerbottommargin=10pt]
      {\centering
      \textbf{\sffamily\small TECHNOLOGY STACK}\\[6pt]
      \small
      \textbf{Frontend:}\ SolidJS 1.9, TailwindCSS 4, @solidjs/router, Kobalte UI\\[2pt]
      \textbf{Backend:}\ Rust 2021, Tokio, rusqlite (bundled), Axum, parking\_lot\\[2pt]
      \textbf{Bridge:}\ Tauri v2 IPC (invoke/emit), Tauri Plugin ecosystem\\[2pt]
      \textbf{Toolchain:}\ yt-dlp, FFmpeg, aria2c, Bun, Vite, GitHub Actions
      \par}
    \end{mdframed}

    \vspace{1.8cm}
    {\small\color{MidGray}
      This document describes the complete systems architecture, engineering decisions,\\
      data-flow design, and concurrency models of the Synclime desktop application.
    }
  \end{center}
  \vfill
  \begin{center}
    \textcolor{LightGray}{\footnotesize
      Produced from a complete file-by-file reading of every source file in the repository.\\
      All architectural details reflect the actual implementation.
    }
  \end{center}
\end{titlepage}

\tableofcontents
\newpage

% ═══════════════════════════════════════════════════════════════════════════════
% DEDICATED FULL-PAGE ARCHITECTURE DIAGRAM
% ═══════════════════════════════════════════════════════════════════════════════
\section{Complete System Architecture Diagram}

The following diagram presents the full component hierarchy of Synclime, from the
smallest individual modules up through layers of increasing abstraction to the
top-level runtime boundary. Each colour identifies a distinct architectural tier.
Arrows denote data flow direction; double-headed arrows indicate bidirectional
communication.

% \clearpage
\thispagestyle{fancy}

\begin{figure}[p]
\centering
\resizebox{\textwidth}{!}{%
\begin{tikzpicture}[
  remember picture,
  node distance=0.35cm and 0.3cm,
  % Styles
  atom/.style={rectangle, rounded corners=2pt, minimum width=2.0cm, minimum height=0.55cm,
               draw, thick, align=center, font=\tiny\sffamily},
  module/.style={rectangle, rounded corners=3pt, minimum width=2.3cm, minimum height=0.7cm,
                 draw, thick, align=center, font=\scriptsize\sffamily\bfseries},
  layer/.style={rectangle, rounded corners=5pt, draw, thick, inner sep=5pt},
  ipc/.style={rectangle, rounded corners=3pt, minimum width=10cm, minimum height=0.75cm,
              draw, thick, align=center, font=\small\sffamily\bfseries},
  arr/.style={-Stealth, thick},
  biarr/.style={Stealth-Stealth, thick},
  groupLabel/.style={font=\scriptsize\sffamily\bfseries},
]

% ══════════════════════════════════════════════════════
% ROW 0: Individual frontend atoms
% ══════════════════════════════════════════════════════

% -- SolidJS Routes atoms --
\node[atom, fill=DiagFELight, draw=DiagFE] (r1) at (0,0)        {Home / Inbox};
\node[atom, fill=DiagFELight, draw=DiagFE, right=0.2cm of r1] (r2) {Queue / Library};
\node[atom, fill=DiagFELight, draw=DiagFE, right=0.2cm of r2] (r3) {Detail Views};
\node[atom, fill=DiagFELight, draw=DiagFE, right=0.2cm of r3] (r4) {Sites / Logs};
\node[atom, fill=DiagFELight, draw=DiagFE, right=0.2cm of r4] (r5) {Settings / About};

% -- Store atoms (right cluster) --
\node[atom, fill=DiagFELight, draw=DiagFE, right=0.8cm of r5]   (s1) {useQueueStore};
\node[atom, fill=DiagFELight, draw=DiagFE, right=0.2cm of s1]  (s2) {useParseStore};
\node[atom, fill=DiagFELight, draw=DiagFE, right=0.2cm of s2]  (s3) {useUIStore};

% ── Module groupings for frontend ──
% Routes module
\begin{scope}[on background layer]
  \node[layer, fill=DiagFELight!40, draw=DiagFE,
        fit=(r1)(r2)(r3)(r4)(r5),
        label={[groupLabel, color=DiagFE]above:12 SolidJS Route Pages}] (routes) {};
\end{scope}

% Stores module
\begin{scope}[on background layer]
  \node[layer, fill=DiagFELight!40, draw=DiagFE,
        fit=(s1)(s2)(s3),
        label={[groupLabel, color=DiagFE]above:Reactive Stores}] (stores) {};
\end{scope}

% Layout node (component, to the right of stores)
\node[module, fill=DiagFELight, draw=DiagFE, right=0.6cm of s3, yshift=0cm] (layout)
  {MainLayout\\+ Sidebar\\+ TitleBar};

% Storage adapter
\node[atom, fill=DiagFELight, draw=DiagFE, below=0.2cm of s2] (adapter)
  {nativeStorageAdapter};

% ══════════ FRONTEND LAYER wrapper ══════════
\begin{scope}[on background layer]
  \node[layer, fill=DiagFELight!20, draw=DiagFE, very thick, inner sep=10pt,
        fit=(routes)(stores)(layout)(adapter),
        label={[font=\normalsize\sffamily\bfseries, color=DiagFE]above:FRONTEND LAYER — SolidJS 1.9 WebView (Tauri)}] (feLayer) {};
\end{scope}

% ══════════════════════════════════════════════════════
% IPC BRIDGE
% ══════════════════════════════════════════════════════
\node[ipc, fill=DiagIPCLight, draw=DiagIPC, below=0.6cm of feLayer.south]
  (ipc) {TAURI v2 IPC BRIDGE \quad
         \texttt{invoke()} $\xrightarrow{\text{typed JSON}}$ Rust \qquad
         Rust $\xrightarrow{\text{typed JSON}}$ \texttt{emit()} \quad
         \textit{\tiny (process boundary — no shared memory)}};

% ══════════════════════════════════════════════════════
% ROW: Command module atoms
% ══════════════════════════════════════════════════════
\node[atom, fill=DiagBELight, draw=DiagBE, below=0.6cm of ipc.south, xshift=-4.0cm]
  (cq) {queue.rs\\11 cmds};
\node[atom, fill=DiagBELight, draw=DiagBE, right=0.2cm of cq]
  (cc) {config.rs\\13 cmds};
\node[atom, fill=DiagBELight, draw=DiagBE, right=0.2cm of cc]
  (cd) {discovery.rs\\2 cmds};
\node[atom, fill=DiagBELight, draw=DiagBE, right=0.2cm of cd]
  (ci) {inbox.rs\\8 cmds};
\node[atom, fill=DiagBELight, draw=DiagBE, right=0.2cm of ci]
  (cl) {logs.rs\\5 cmds};
\node[atom, fill=DiagBELight, draw=DiagBE, right=0.2cm of cl]
  (cb) {clipboard.rs\\1 cmd};

\begin{scope}[on background layer]
  \node[layer, fill=DiagBELight!40, draw=DiagBE,
        fit=(cq)(cc)(cd)(ci)(cl)(cb),
        label={[groupLabel, color=DiagBE]above:Command Handlers — 30+ typed IPC commands across 6 modules}]
    (cmdGroup) {};
\end{scope}

% ══════════════════════════════════════════════════════
% ROW: AppEngineState
% ══════════════════════════════════════════════════════
\node[module, fill=DiagBELight, draw=DiagBE, below=0.4cm of cmdGroup.south, xshift=-2.5cm]
  (state) {AppEngineState\\{\tiny Arc<parking\_lot::*>}};

% State sub-components
\node[atom, fill=DiagBELight, draw=DiagBE, right=0.3cm of state] (sem)
  {Semaphore Pool\\{\tiny RwLock<Arc<Sem>>}};
\node[atom, fill=DiagBELight, draw=DiagBE, right=0.3cm of sem] (sig)
  {Signal Channel\\{\tiny mpsc::Sender}};
\node[atom, fill=DiagBELight, draw=DiagBE, right=0.3cm of sig] (cache)
  {Progress Cache\\{\tiny Mutex<HashMap>}};
\node[atom, fill=DiagBELight, draw=DiagBE, right=0.3cm of cache] (preg)
  {Process Registry\\{\tiny RwLock<HashMap>}};

\begin{scope}[on background layer]
  \node[layer, fill=DiagBELight!40, draw=DiagBE,
        fit=(state)(sem)(sig)(cache)(preg),
        label={[groupLabel, color=DiagBE]above:AppEngineState — Shared Arc (Tauri managed)}]
    (stateGroup) {};
\end{scope}

% ══════════════════════════════════════════════════════
% ROW: Engine atoms
% ══════════════════════════════════════════════════════
\node[atom, fill=DiagBELight, draw=DiagBE, below=0.4cm of stateGroup.south, xshift=-3.2cm]
  (ew) {Download\\Worker Task};
\node[atom, fill=DiagBELight, draw=DiagBE, right=0.3cm of ew]
  (pp) {Dual Progress\\Parser};
\node[atom, fill=DiagBELight, draw=DiagBE, right=0.3cm of pp]
  (cpb) {yt-dlp CMD\\Builder};
\node[atom, fill=DiagBELight, draw=DiagBE, right=0.3cm of cpb]
  (ck) {Cookie RAII\\Guard};
\node[atom, fill=DiagBELight, draw=DiagBE, right=0.3cm of ck]
  (cw) {Cancellation\\Worker};
\node[atom, fill=DiagBELight, draw=DiagBE, right=0.3cm of cw]
  (fw) {1-sec Flush\\Worker};

\begin{scope}[on background layer]
  \node[layer, fill=DiagBELight!40, draw=DiagBE,
        fit=(ew)(pp)(cpb)(ck)(cw)(fw),
        label={[groupLabel, color=DiagBE]above:engine/process.rs — Download Pipeline (748 lines)}]
    (engineGroup) {};
\end{scope}

% Axum server (right of engine group)
\node[module, fill=DiagAxumLight, draw=DiagAxum, right=0.5cm of engineGroup.east]
  (axum) {Axum 0.7\\API Server\\{\tiny :14221--14230}};
\node[atom, fill=DiagAxumLight, draw=DiagAxum, below=0.2cm of axum] (axumR1)
  {GET /health};
\node[atom, fill=DiagAxumLight, draw=DiagAxum, below=0.2cm of axumR1] (axumR2)
  {POST /add};

% ══════════ BACKEND LAYER wrapper ══════════
\begin{scope}[on background layer]
  \node[layer, fill=DiagBELight!10, draw=DiagBE, very thick, inner sep=12pt,
        fit=(cmdGroup)(stateGroup)(engineGroup)(axum)(axumR1)(axumR2),
        label={[font=\normalsize\sffamily\bfseries, color=DiagBE]below:NATIVE RUST BACKEND — Tokio Async Runtime (Rust 2021)}]
    (beLayer) {};
\end{scope}

% ══════════════════════════════════════════════════════
% PERSISTENCE ROW
% ══════════════════════════════════════════════════════
\node[module, fill=DiagDBLight, draw=DiagDB, below=0.6cm of beLayer.south, xshift=-2.5cm]
  (sqlite) {SQLite\\synclime\_core.db\\{\tiny 9 tables, 4 indexes}};

% DB tables atom cluster
\node[atom, fill=DiagDBLight, draw=DiagDB, right=0.3cm of sqlite] (t1)
  {download\_jobs\\parsed\_files};
\node[atom, fill=DiagDBLight, draw=DiagDB, right=0.25cm of t1] (t2)
  {cookie\_profiles\\proxy\_profiles};
\node[atom, fill=DiagDBLight, draw=DiagDB, right=0.25cm of t2] (t3)
  {site\_configs\\app\_settings};
\node[atom, fill=DiagDBLight, draw=DiagDB, right=0.25cm of t3] (t4)
  {inbox\_urls\\parse\_logs};
\node[atom, fill=DiagDBLight, draw=DiagDB, right=0.25cm of t4] (t5)
  {error\_logs\\app\_fallback};

\begin{scope}[on background layer]
  \node[layer, fill=DiagDBLight!40, draw=DiagDB,
        fit=(sqlite)(t1)(t2)(t3)(t4)(t5),
        label={[groupLabel, color=DiagDB]below:PERSISTENCE — rusqlite bundled (SQLite amalgamation compiled-in)}]
    (dbLayer) {};
\end{scope}

% Subprocess node
\node[module, fill=DiagProcLight, draw=DiagProc, right=0.6cm of dbLayer.east]
  (ytdlp) {yt-dlp\\+ aria2c\\Subprocess};
\node[atom, fill=DiagProcLight, draw=DiagProc, below=0.2cm of ytdlp] (ffmpeg)
  {FFmpeg\\muxer};

\begin{scope}[on background layer]
  \node[layer, fill=DiagProcLight!30, draw=DiagProc, inner sep=6pt,
        fit=(ytdlp)(ffmpeg),
        label={[groupLabel, color=DiagProc]below:OS Processes}]
    (procLayer) {};
\end{scope}

% ══════════════════════════════════════════════════════
% ARROWS
% ══════════════════════════════════════════════════════

% Frontend internals
\draw[arr, DiagFE] (routes.east) -- (stores.west);
\draw[arr, DiagFE] (adapter.west) -| (stores.south);

% FE <-> IPC
\draw[biarr, DiagIPC, very thick] (feLayer.south) -- (ipc.north);

% IPC <-> Commands
\draw[biarr, DiagIPC, very thick] (ipc.south) -- (cmdGroup.north);

% Commands <-> State
\draw[biarr, DiagBE] (cmdGroup.south) -- (stateGroup.north);

% State -> Engine
\draw[arr, DiagBE] (stateGroup.south) -- (engineGroup.north);

% Engine -> Axum (horizontal)
\draw[arr, DiagAxum] (engineGroup.east) -- (axum.west);

% State <-> DB
\draw[biarr, DiagDB, thick] (stateGroup.south) -- ++(0,-0.3) -| (dbLayer.north);

% Cmd <-> DB
\draw[biarr, DiagDB, thick] (cmdGroup.south) -- ++(0,-1.2) -| (dbLayer.north);

% Engine -> Process
\draw[biarr, DiagProc, thick] (engineGroup.south) -- ++(0,-0.3) -| (procLayer.north);

% Axum -> DB
\draw[arr, DiagDB] (axumR2.south) -- ++(0,-0.3) -| (dbLayer.north);

\end{tikzpicture}
}%
\caption{%
  \textbf{Complete Synclime Component Architecture.}
  Colour key: \textcolor{DiagFE}{$\blacksquare$} Frontend (SolidJS),
              \textcolor{DiagIPC}{$\blacksquare$} IPC Bridge (Tauri),
              \textcolor{DiagBE}{$\blacksquare$} Rust Backend (Tokio),
              \textcolor{DiagDB}{$\blacksquare$} Persistence (SQLite),
              \textcolor{DiagProc}{$\blacksquare$} OS Subprocesses,
              \textcolor{DiagAxum}{$\blacksquare$} Embedded API (Axum).
}
\end{figure}

\clearpage

% ═══════════════════════════════════════════════════════════════════════════════
\section{Executive Overview}
% ═══════════════════════════════════════════════════════════════════════════════

Synclime is a production-grade cross-platform native desktop application engineered to solve
a concrete, non-trivial systems problem: orchestrating high-volume media downloads from
the internet without freezing the user interface, corrupting shared state, or leaking
OS-level processes.

The application ingests video, audio, subtitle, and direct document URLs; resolves metadata
via an integrated \texttt{yt-dlp} subprocess pipeline; persists download job configurations
in a relational SQLite database with enforced foreign-key integrity; and orchestrates
concurrent download workers through a Tokio-backed async runtime with a bounded semaphore
pool. Real-time progress telemetry is debounced through an in-memory progress cache, flushed
to both the database and the UI at 1-second intervals---a deliberate design trade-off that
eliminates the most common cause of UI thread freezes in download-manager applications.

\begin{insightbox}{Central Design Constraint}
The entire architecture flows from one core constraint: \emph{the UI must never block on I/O,
the database must never absorb unbounded write storms, and OS-level processes must be safely
killable at any moment.} Every structural decision documented below is a direct consequence
of that constraint.
\end{insightbox}

The application ships on five binary targets (Windows x64/ARM64, macOS ARM64, Linux
x64/ARM64) assembled automatically through a GitHub Actions matrix pipeline with
platform-specific native dependency management.

\subsection{Key Capabilities}

\begin{itemize}[leftmargin=1.6em, itemsep=3pt]
  \item \textbf{Concurrent download queue:} Semaphore-bounded worker pool, dynamically
        reconfigurable at runtime without restarting any running jobs.
  \item \textbf{Debounced progress telemetry:} In-memory HashMap absorbs thousands of
        yt-dlp stdout lines per second; a background Tokio task flushes to SQLite and
        emits one Tauri event per job per second.
  \item \textbf{Safe process cancellation:} Pause signals route through an mpsc channel
        to a dedicated cancellation worker; on Unix, whole process groups are killed
        (\texttt{kill -9 -<pgid>}) to eliminate aria2c zombie children.
  \item \textbf{Cookie-secured downloads:} Netscape-format session cookies are written
        to \texttt{0o600}-permission temporary files and deleted by a Rust \texttt{Drop}
        guard after each job, regardless of outcome.
  \item \textbf{Browser extension inbox:} An embedded Axum HTTP server listens on
        \texttt{127.0.0.1:14221--14230} and persists URLs from companion extensions
        directly to SQLite with real-time badge updates on the UI.
  \item \textbf{Domain routing rules:} Per-domain site configurations associate cookie
        vaults and SOCKS5/HTTP proxy profiles via SQLite foreign keys, with live
        validation badges detecting broken credential references.
  \item \textbf{Dual-format progress parsing:} ANSI-stripped yt-dlp lines and
        aria2c parenthetical percentages are decoded by hand-written character-level
        scanners---not regex---for robustness against format variations.
  \item \textbf{Full audit logging:} Parse operations and download errors are persisted
        to \texttt{parse\_logs} and \texttt{error\_logs} with executed commands, exit
        codes, byte counts, and RFC 3339 timestamps.
\end{itemize}

% ═══════════════════════════════════════════════════════════════════════════════
\section{Rust Backend Design}
% ═══════════════════════════════════════════════════════════════════════════════

\subsection{Central State Object: AppEngineState}

All shared mutable application state is consolidated into one \texttt{AppEngineState} struct
managed by Tauri's \texttt{app.manage()} API. Every command handler receives it as an
injected \texttt{State<'\_, AppEngineState>}---zero-copy access without constructor injection
boilerplate.

\begin{lstlisting}[caption={AppEngineState (lib.rs)}]
pub struct AppEngineState {
    pub pool_semaphore: Arc<parking_lot::RwLock<Arc<Semaphore>>>,
    pub db_path: std::path::PathBuf,
    pub db_conn: Arc<parking_lot::Mutex<Connection>>,
    pub active_processes: Arc<ActiveProcessRegistry>,
    pub signal_tx: mpsc::Sender<QueueSignal>,
    pub progress_cache: Arc<parking_lot::Mutex<HashMap<String, ProgressSnapshot>>>,
}
\end{lstlisting}

\begin{insightbox}{Why parking\_lot over std::sync?}
\texttt{parking\_lot::Mutex} and \texttt{RwLock} are smaller, faster under contention,
and lack the poisoning semantics of the standard library (panics in a holding thread do
not corrupt shared state for other threads). For a data path hit once per second by the
flush worker and on-demand by every command handler, this is a deliberate performance choice.
\end{insightbox}

The \texttt{pool\_semaphore} is wrapped as \texttt{RwLock<Arc<Semaphore>>} rather than a
plain \texttt{Semaphore}. This enables the concurrency limit to be hot-swapped at runtime:
the \texttt{update\_concurrency\_limit} command writes a new \texttt{Arc<Semaphore>} under
the write lock without invalidating any currently-held permits in running workers.

\subsection{Concurrency Model: The Semaphore Pool}

Download workers are dispatched as detached Tokio tasks:

\begin{lstlisting}[caption={Non-blocking task dispatch (commands/queue.rs)}]
tauri::async_runtime::spawn(async move {
    let _ = execute_download_worker(app_handle, worker_slug).await;
});
\end{lstlisting}

Inside \texttt{execute\_download\_worker}, the worker acquires a semaphore permit before
spawning the OS process. The permit is held for the full lifetime of the download---released
when the function returns. No more than $N$ concurrent yt-dlp processes can exist, where $N$
is the user-configured concurrency limit. The owned permit pattern guarantees release
through RAII even on panics or early returns.

\subsection{The 1-Second Progress Flush Architecture}

A naive implementation would write every yt-dlp stdout line directly to SQLite and emit a
Tauri event per line---thousands of writes per second per download, which reliably causes
UI freezes and database lock contention.

Synclime resolves this with a two-stage pipeline:

\begin{enumerate}[leftmargin=2em]
  \item \textbf{Stage 1 (per stdout line):} Progress data is written to a
        \texttt{HashMap<String, ProgressSnapshot>} under a Mutex---a sub-microsecond
        in-memory operation.
  \item \textbf{Stage 2 (once per second):} A background Tokio task wakes, locks the
        cache, bulk-writes all snapshots to SQLite, emits one Tauri event per active job,
        then clears the cache.
\end{enumerate}

\begin{lstlisting}[caption={1-second flush worker (lib.rs)}]
tauri::async_runtime::spawn(async move {
    loop {
        tokio::time::sleep(Duration::from_secs(1)).await;
        let mut cache = flush_cache.lock();
        if !cache.is_empty() {
            let conn = flush_conn.lock();
            for (slug, snap) in cache.iter() {
                let _ = conn.execute(
                    "UPDATE download_jobs
                     SET progress=?1, tracking_message=?2, status=?3,
                         updated_at=datetime('now')
                     WHERE slug=?4;",
                    params![snap.progress, snap.status_message, snap.status, slug],
                );
                let _ = tic_emitter.emit("download-progress-token",
                    json!({"slug":slug,"progress":snap.progress,
                           "message":&snap.status_message,"status":&snap.status}));
            }
            cache.clear();
        }
    }
});
\end{lstlisting}

\subsection{Safe Process Cancellation via mpsc Channel}

\begin{lstlisting}[caption={Cancellation worker (lib.rs)}]
async fn start_cancellation_worker(
    mut rx: mpsc::Receiver<QueueSignal>,
    registry: Arc<ActiveProcessRegistry>,
) {
    while let Some(QueueSignal::PauseJob(slug)) = rx.recv().await {
        let process = {
            let mut instances = registry.instances.write();
            instances.remove(&slug)
        };
        if let Some(mut child) = process {
            #[cfg(unix)]
            if let Some(pid) = child.id() {
                // Kill entire process GROUP: yt-dlp + any aria2c children
                let _ = Command::new("kill")
                    .args(&["-9", &format!("-{}", pid)])
                    .status();
            }
            #[cfg(not(unix))]
            let _ = child.kill().await;
        }
    }
}
\end{lstlisting}

On Unix, \texttt{kill -9 -<pgid>} terminates the entire process group, preventing aria2c
fragment workers from becoming zombie processes after yt-dlp itself exits.

\subsection{Cookie Security: RAII Drop Guard}

\begin{lstlisting}[caption={Cookie file RAII cleanup (engine/process.rs)}]
struct CookieFileCleanup(Option<std::path::PathBuf>);
impl Drop for CookieFileCleanup {
    fn drop(&mut self) {
        if let Some(ref path) = self.0 {
            let _ = std::fs::remove_file(path); // runs on success, error, or panic
        }
    }
}
let _cleanup_guard = CookieFileCleanup(temp_cookie_path.clone());
\end{lstlisting}

\subsection{Resilient Progress Parsing: Hand-Written Character Scanners}

Rather than brittle regex, Synclime uses two hand-written character-level scanners:

\begin{description}[style=nextline, leftmargin=1.6em]
  \item[\texttt{parse\_progress\_line}] Strips ANSI sequences character-by-character using
        an \texttt{in\_ansi} state machine flag, then backward-scans from the first
        \texttt{\%} to extract the numeric percentage.
  \item[\texttt{parse\_aria2\_progress}] Scans all \texttt{\%} characters, backward-searches
        for the nearest \texttt{(}, and extracts aria2c's \texttt{(99\%)} format.
\end{description}

\begin{lstlisting}[caption={Dual-parser chaining (engine/process.rs)}]
let parsed_progress = parse_progress_line(clean_line)
    .or_else(|| parse_aria2_progress(clean_line));
\end{lstlisting}

Four unit tests in \texttt{\#[cfg(test)] mod tests} cover standard output, whitespace
edge cases, ANSI-prefixed lines, and non-progress lines.

% ═══════════════════════════════════════════════════════════════════════════════
\section{Database Architecture}
% ═══════════════════════════════════════════════════════════════════════════════

\subsection{Schema Design}

\begin{table}[H]
\centering\small
\renewcommand{\arraystretch}{1.35}
\begin{tabularx}{\linewidth}{>{\bfseries\sffamily}lXl}
\toprule
Table & Purpose & Key Relations \\
\midrule
cookie\_profiles & Netscape-format cookie blobs per domain & --- \\
proxy\_profiles & SOCKS5/HTTP proxy endpoint strings & --- \\
site\_configs & Per-domain routing rules & FK $\to$ cookie, proxy \\
app\_settings & KV store: paths, limits, active port & --- \\
parsed\_files & yt-dlp extracted metadata manifests & FK $\to$ site\_configs, self \\
download\_jobs & Queue items with status/progress & FK $\to$ parsed\_files, cookies, proxy, self \\
parse\_logs & Audit log for extraction runs & FK $\to$ parsed\_files (CASCADE) \\
error\_logs & Per-job error messages + commands & FK $\to$ download\_jobs (CASCADE) \\
inbox\_urls & Browser extension ingestion buffer & --- \\
\bottomrule
\end{tabularx}
\caption{Database schema summary (9 tables)}
\end{table}

\subsection{Foreign Key Integrity}

\texttt{PRAGMA foreign\_keys = ON} is executed on every connection open. Referential
actions are tuned per relationship:

\begin{itemize}[leftmargin=1.6em, itemsep=2pt]
  \item \texttt{site\_configs} $\to$ credentials: \texttt{ON DELETE SET NULL}
  \item \texttt{download\_jobs} $\to$ download\_jobs (self): \texttt{ON DELETE CASCADE}
  \item \texttt{parse\_logs} / \texttt{error\_logs}: \texttt{ON DELETE CASCADE}
\end{itemize}

\subsection{Performance Indexes}

\begin{lstlisting}[caption={Partial indexes (lib.rs)}]
CREATE INDEX IF NOT EXISTS idx_parsed_files_parent
  ON parsed_files (parent_playlist_slug)
  WHERE parent_playlist_slug IS NOT NULL;

CREATE INDEX IF NOT EXISTS idx_download_jobs_parsed_file
  ON download_jobs (parsed_file_slug)
  WHERE parsed_file_slug IS NOT NULL;

CREATE INDEX IF NOT EXISTS idx_download_jobs_subtitles
  ON download_jobs (associated_media_job_slug)
  WHERE file_type = 'subtitle' AND associated_media_job_slug IS NOT NULL;

CREATE INDEX IF NOT EXISTS idx_download_jobs_queue_priority
  ON download_jobs (status, priority_index DESC, created_at ASC);
\end{lstlisting}

\subsection{The app\_fallback Sentinel Pattern}

On startup, sentinel rows with slug \texttt{'app\_fallback'} are inserted into both
\texttt{parsed\_files} and \texttt{download\_jobs} via \texttt{INSERT OR IGNORE}. Error
logging commands fall back to this sentinel if the target job has been deleted---preventing
foreign key violations from silently swallowing diagnostic data.

% ═══════════════════════════════════════════════════════════════════════════════
\section{IPC Command Architecture}
% ═══════════════════════════════════════════════════════════════════════════════

\subsection{Command Module Map}

\begin{table}[H]
\centering\small
\renewcommand{\arraystretch}{1.3}
\begin{tabular}{>{\ttfamily}lrl}
\toprule
Module & Cmds & Responsibility \\
\midrule
commands/queue.rs & 11 & Job lifecycle: start, pause, insert, delete, concurrency settings \\
commands/config.rs & 13 & Cookie/proxy profiles, site configs, download path \\
commands/discovery.rs & 2 & yt-dlp metadata extraction, parsed file persistence \\
commands/inbox.rs & 8 & Inbox CRUD, remote update fetch, API port query \\
commands/logs.rs & 5 & Error/parse log retrieval, insertion, bulk clear \\
commands/clipboard.rs & 1 & URL validation + UTM tracking parameter stripping \\
\bottomrule
\end{tabular}
\caption{Tauri IPC command module map}
\end{table}

\subsection{Job Configuration Resolution: Multi-Level SQL JOIN}

\begin{lstlisting}[caption={Multi-join config resolution (engine/process.rs)}]
SELECT
    COALESCE(j.direct_url, p.url) as target_url,
    COALESCE(sc_cookie.cookie_data, c.cookie_data) as cookie_data,
    COALESCE(sc_proxy.proxy_string, pr.proxy_string) as proxy_string, ...
FROM download_jobs j
LEFT JOIN parsed_files p ON j.parsed_file_slug = p.slug
LEFT JOIN parsed_files parent_p ON p.parent_playlist_slug = parent_p.slug
LEFT JOIN site_configs sc ON
    COALESCE(p.site_config_slug, parent_p.site_config_slug) = sc.slug
LEFT JOIN cookie_profiles sc_cookie ON sc.cookie_profile_slug = sc_cookie.slug
LEFT JOIN proxy_profiles sc_proxy ON sc.proxy_profile_slug = sc_proxy.slug
LEFT JOIN cookie_profiles c ON j.cookie_profile_slug = c.slug
LEFT JOIN proxy_profiles pr ON j.proxy_profile_slug = pr.slug
WHERE j.slug = ?1;
\end{lstlisting}

Priority: \textbf{1)} Direct per-job override; \textbf{2)} Site config on the parsed file;
\textbf{3)} Site config inherited from parent playlist. Resolved in a single round-trip.

% ═══════════════════════════════════════════════════════════════════════════════
\section{Frontend Architecture}
% ═══════════════════════════════════════════════════════════════════════════════

\subsection{SolidJS Reactive Stores}

\begin{table}[H]
\centering\small
\renewcommand{\arraystretch}{1.3}
\begin{tabular}{>{\ttfamily}lll}
\toprule
Store & State Shape & Persistence \\
\midrule
useQueueStore & DownloadJob[], progress map & In-memory; refreshed from SQLite \\
useParseStore & ParsedFile[], isParsing & Tauri Store plugin (binary file) \\
useUIStore & theme, path, badges, sidebar & Tauri Store plugin (binary file) \\
\bottomrule
\end{tabular}
\caption{Frontend reactive stores}
\end{table}

\subsection{Real-Time Event Architecture}

\begin{lstlisting}[caption={Global event listener (store/useQueueStore.ts)}]
import("@tauri-apps/api/event").then(({ listen }) => {
  listen("download-progress-token", (event: any) => {
    const { slug, progress, message, status } = event.payload;
    useQueueStore.updateJobProgress(slug, progress, message);
    if (status) useQueueStore.updateJobStatus(slug, status as any);
  });
});
\end{lstlisting}

\subsection{Custom TitleBar and Window Safety}

The Tauri window uses \texttt{"decorations": false, "transparent": true}. The \emph{Quit App}
button triggers a native OS \texttt{ask()} dialog warning active downloads will be
interrupted. F5, Ctrl+R, and DevTools shortcuts are intercepted and blocked at the
\texttt{keydown} level to prevent accidental WebView state destruction.

% ═══════════════════════════════════════════════════════════════════════════════
\section{Embedded Axum API Server}
% ═══════════════════════════════════════════════════════════════════════════════

\begin{table}[H]
\centering\small
\renewcommand{\arraystretch}{1.3}
\begin{tabular}{llll}
\toprule
Endpoint & Method & Input & Effect \\
\midrule
\texttt{/health} & GET & --- & JSON status, version \\
\texttt{/add} & POST & \texttt{\{"url": "..."\}} & INSERT to inbox\_urls; emit \texttt{inbox-updated} \\
\bottomrule
\end{tabular}
\caption{Axum local API endpoints}
\end{table}

Port range \texttt{14221--14230} is scanned for the first available binding; the active
port is persisted to \texttt{app\_settings} for frontend discovery. The \texttt{/add}
handler shares the same \texttt{Arc<Mutex<Connection>>} as the rest of the application,
preserving the single-writer SQLite model.

% ═══════════════════════════════════════════════════════════════════════════════
\section{CI/CD Pipeline and Cross-Platform Build Matrix}
% ═══════════════════════════════════════════════════════════════════════════════

\begin{table}[H]
\centering\small
\renewcommand{\arraystretch}{1.3}
\begin{tabular}{lll}
\toprule
Runner & Rust Target & Notes \\
\midrule
macos-latest & aarch64-apple-darwin & Apple Silicon \\
ubuntu-22.04 & x86\_64-unknown-linux-gnu & x64 Linux \\
ubuntu-22.04-arm & aarch64-unknown-linux-gnu & ARM64 Linux \\
windows-latest & x86\_64-pc-windows-msvc & Windows x64 \\
windows-latest & aarch64-pc-windows-msvc & Windows ARM64 \\
\bottomrule
\end{tabular}
\caption{GitHub Actions 5-target build matrix}
\end{table}

\begin{itemize}[leftmargin=1.6em, itemsep=2pt]
  \item \textbf{Linux:} apt installs WebKit2GTK, HarfBuzz, Graphite2, ICU; \texttt{RUSTFLAGS}
        force-links \texttt{-lstdc++} and \texttt{-lgraphite2}.
  \item \textbf{macOS:} Homebrew installs icu4c and libpng; \texttt{PKG\_CONFIG\_PATH}
        covers both Intel and Apple Silicon prefix locations.
  \item \textbf{Windows:} vcpkg installs ICU, HarfBuzz, Graphite2, fontconfig with
        architecture-aware triplets; \texttt{CXXFLAGS=/std:c++17} set for MSVC.
\end{itemize}

% ═══════════════════════════════════════════════════════════════════════════════
\section{Security Model}
% ═══════════════════════════════════════════════════════════════════════════════

\subsection{Tauri Capability Model (Least Privilege)}

\texttt{capabilities/default.json} explicitly enumerates the permissions granted to the
main window. No filesystem read/write capability is granted to the WebView; all file
access is mediated through typed Rust commands.

\subsection{Cookie File Security}

\begin{itemize}[leftmargin=1.6em, itemsep=2pt]
  \item Unix: \texttt{0o600} (owner-read-write only) via \texttt{OpenOptionsExt}
  \item RAII \texttt{Drop} guard deletes the file after the download exits---success, error, or panic
  \item Startup sweep removes any \texttt{synclime\_cookie\_*.txt} leftover from crashes
  \item Nanosecond-timestamped filenames prevent concurrent-download collisions
\end{itemize}

% ═══════════════════════════════════════════════════════════════════════════════
\section{Engineering Decisions Summary}
% ═══════════════════════════════════════════════════════════════════════════════

\begin{techbox}{Key Technical Choices and Their Rationale}
\begin{description}[style=nextline, leftmargin=1.0em, itemsep=5pt]
  \item[\textbf{SolidJS over React}]
    Fine-grained reactivity: only the changed signal re-renders. For a queue with
    sub-second progress updates, this is measurably better than Virtual DOM diffing.
  \item[\textbf{rusqlite bundled feature}]
    SQLite amalgamation compiled into the binary: zero runtime dependency on system
    SQLite, identical behaviour across all five target platforms.
  \item[\textbf{parking\_lot over std::sync}]
    Lower lock overhead under contention; no mutex poisoning semantics.
  \item[\textbf{mpsc channel for cancellation}]
    Decouples the IPC command handler (must return fast) from blocking OS kill calls.
  \item[\textbf{Arc<RwLock<Arc<Semaphore>>>}]
    Enables atomic hot-swap of the semaphore pool when the concurrency limit changes,
    without disrupting currently-acquiring workers.
  \item[\textbf{Character-level progress scanners}]
    Immune to ANSI escape sequence variants and Unicode edge cases that would break regex.
  \item[\textbf{Bun + frozen lockfile}]
    Deterministic CI builds; Bun's native TypeScript execution speeds the development loop.
  \item[\textbf{Port range scanning 14221--14230}]
    Handles port conflicts gracefully without user configuration.
\end{description}
\end{techbox}

% ═══════════════════════════════════════════════════════════════════════════════
\section{Codebase Metrics}
% ═══════════════════════════════════════════════════════════════════════════════

\begin{table}[H]
\centering\small
\renewcommand{\arraystretch}{1.3}
\begin{tabular}{lrrl}
\toprule
\textbf{Layer} & \textbf{Files} & \textbf{Approx. Lines} & \textbf{Language} \\
\midrule
Rust backend (lib.rs + 5 cmd modules + engine) & 11 & $\sim$2,600 & Rust \\
SolidJS routes (12 page components) & 12 & $\sim$3,600 & TypeScript/TSX \\
Layouts \& shared components & 4 & $\sim$700 & TypeScript/TSX \\
Reactive stores \& storage adapter & 4 & $\sim$280 & TypeScript \\
Core types \& logger & 4 & $\sim$280 & TypeScript \\
Configuration \& CI & 6 & $\sim$350 & --- \\
\midrule
\textbf{Total} & \textbf{41} & \textbf{$\sim$7,800} & \\
\bottomrule
\end{tabular}
\caption{Codebase metrics by architectural layer}
\end{table}

% ═══════════════════════════════════════════════════════════════════════════════
\section{Conclusion}
% ═══════════════════════════════════════════════════════════════════════════════

Synclime demonstrates a coherent set of senior engineering practices applied consistently
across a full-stack native desktop application:

\begin{itemize}[leftmargin=1.6em, itemsep=4pt]
  \item \textbf{Concurrent systems design:} Semaphore-bounded pools, mpsc cancellation
        channels, and debounced flush architecture prevent UI freeze, ghost processes,
        and database write storms.
  \item \textbf{Relational data modelling:} Normalised schema with foreign key integrity,
        partial indexes, carefully chosen referential actions, and sentinel pattern for
        orphan-safe audit logging.
  \item \textbf{Security by design:} Least-privilege Tauri capabilities, RAII-guarded
        temporary secrets, process group termination on Unix, startup artifact cleanup.
  \item \textbf{Resilient parsing:} Hand-written character scanners over fragile regex
        for progress output subject to ANSI codes and format variations.
  \item \textbf{Cross-platform delivery:} 5-target CI/CD matrix with platform-specific
        native dependency management producing native installers automatically.
  \item \textbf{Integration depth:} Embedded Axum server, reactive event pipeline, and
        dual-backend storage adapter create a seamless flow from browser capture to
        download completion with no polling at any layer.
\end{itemize}

\begin{insightbox}{Architecture Quality Signal}
The codebase is free of several patterns diagnostic of intermediate-level work: no global
mutable statics, no \texttt{unwrap()} in production paths (errors propagate via
\texttt{Result<T, String>}), no raw thread spawning (all concurrency through Tokio), no
shared state without synchronisation primitives, and no sleep-loop polling for coordination.
The architecture is \emph{reactive at both ends}: Rust pushes typed events; SolidJS
consumes them through a module-level listener registered once. Neither side polls.

The result is a production-grade native desktop application that handles real concurrency
and real failure modes with principled, deliberate engineering.
\end{insightbox}

\vspace{1.5cm}
\noindent{\color{RuleGray}\rule{\linewidth}{0.4pt}}
\vspace{0.3cm}

\noindent\small\color{MidGray}
\textbf{Synclime Technical Architecture Report} \textemdash{} v1.0.0 \quad
Author: Ahmed Trooper \quad
Repository: github.com/AhmedTrooper/Synclime\\
License: MIT / Apache 2.0 (Dual) \quad
Generated: July 2026

\end{document}
